% Компиляция: XeLaTeX, два запуска.
% 11 слайдов, включая подробные таблицы результатов.
% Источники чисел: финальные таблицы из run_summaries.
% Комментарии перед слайдами — ориентиры для устного доклада.

\documentclass[aspectratio=169,10pt]{beamer}

\usepackage{fontspec}
\usepackage[russian]{babel}
\setsansfont{Arial}
\setmainfont{Arial}
\setmonofont{Consolas}[Scale=0.85]
\usefonttheme{professionalfonts}

\usepackage{amsmath,amssymb}
\usepackage{booktabs,tabularx,array,colortbl}

% -------------------- Оформление --------------------

\definecolor{Ink}{HTML}{17263C}
\definecolor{Accent}{HTML}{176B75}
\definecolor{Muted}{HTML}{566271}
\definecolor{Tint}{HTML}{EDF4F5}

\setbeamercolor{normal text}{fg=Ink,bg=white}
\setbeamercolor{structure}{fg=Accent}
\setbeamercolor{frametitle}{fg=Ink,bg=white}
\setbeamercolor{block title}{fg=Accent,bg=white}
\setbeamercolor{block body}{fg=Ink,bg=white}
\setbeamercolor{takeaway}{fg=Ink,bg=Tint}
\setbeamercolor{footline}{fg=Muted,bg=white}

\setbeamerfont{frametitle}{
  size=\large,
  series=\bfseries
}
\setbeamerfont{block title}{
  size=\normalsize,
  series=\bfseries
}
\setbeamerfont{footline}{size=\scriptsize}

\setbeamersize{
  text margin left=7mm,
  text margin right=7mm
}
\setbeamertemplate{navigation symbols}{}
\setbeamertemplate{itemize item}{\small\textbullet}

\setbeamertemplate{footline}{%
  \leavevmode
  \begin{beamercolorbox}[
    wd=\paperwidth,
    ht=2.8ex,
    dp=1.4ex,
    leftskip=7mm,
    rightskip=7mm
  ]{footline}
    Квантизация SD 1.5
    \hfill \insertframenumber/\inserttotalframenumber
  \end{beamercolorbox}%
}

\newcommand{\keypoint}[1]{%
  \begin{beamercolorbox}[
    sep=1.8mm,
    wd=\linewidth
  ]{takeaway}
    \small\bfseries #1
  \end{beamercolorbox}%
}

\newcommand{\muted}[1]{{\color{Muted}#1}}
\newcommand{\midblock}{\texttt{mid\_block}}
\newcommand{\rvq}{\textbf{RVQ-AW}}
\newcommand{\score}{\operatorname{Score}}
\newcolumntype{Y}{>{\raggedright\arraybackslash}X}

\title{Квантизация диффузионных моделей машинного обучения}
\subtitle{Activation-aware RVQ для Stable Diffusion 1.5}
\author{Подъяблонский Егор Петрович}
\institute{ВМК МГУ}
\date{2026}

\begin{document}

% ===================================================
% 1. ~20 секунд.
% Представить тему и исследовательский вопрос.
% ===================================================

\begin{frame}[plain]
  \vspace{3mm}

  {\color{Accent}\small
    КУРСОВАЯ РАБОТА
    \quad / \quad
    ПРОМЕЖУТОЧНЫЕ РЕЗУЛЬТАТЫ
    \par
  }

  \vspace{7mm}

  {\LARGE\bfseries
    Квантизация диффузионных\par
    моделей машинного обучения\par
  }

  \vspace{4mm}

  {\large
    Activation-aware RVQ для Stable Diffusion 1.5
    \par
  }

  \vfill

  {\normalsize
    Подъяблонский Егор Петрович\par
  }

  \vspace{2mm}

  {\small\color{Muted}
    ВМК МГУ \quad | \quad 2026\par
  }

  \vspace{3mm}
\end{frame}

% ===================================================
% 2. ~50 секунд.
% Разделить общую задачу сжатия и текущий масштаб проверки.
% ===================================================

\begin{frame}[t]{Задача и границы исследования}
  Диффузионные модели требуют значительных ресурсов.
  В работе исследуется
  \textbf{post-training quantization (PTQ)} весов.

  \vspace{1mm}

  \begin{columns}[T,onlytextwidth]

    \begin{column}{0.47\textwidth}
      \begin{block}{Исследовательская цель}
        Сохранить большую часть качества сильного векторного
        бейзлайна \textbf{VQDM}, но уменьшить время построения
        квантизованного представления.
      \end{block}

      \small
      Узкое место VQDM --- совместный поиск комбинаций
      кодовых слов с помощью \textbf{beam search}.
    \end{column}

    \begin{column}{0.49\textwidth}
      \begin{block}{Что квантуется сейчас}
        \begin{itemize}
          \setlength{\itemsep}{2mm}

          \item U-Net модели
            \textbf{Stable Diffusion 1.5}.

          \item \midblock:
            \textbf{18 Linear/Conv2d-слоёв};
            также проверяется набор из 16 слоёв
            без двух \texttt{time\_emb\_proj}.

          \item \textbf{Только веса, без fine-tuning.}
        \end{itemize}
      \end{block}
    \end{column}

  \end{columns}

  \vfill

  \keypoint{
    Представленные далее результаты --- для 18 слоёв
    \midblock, а не для полного U-Net.
  }
\end{frame}

% ===================================================
% 3. ~45 секунд.
% Не пересказывать статьи: объяснить роль групп методов.
% ===================================================

\begin{frame}[t]{С чем сравнивается предлагаемый метод}
  \begin{columns}[T,onlytextwidth]

    \begin{column}{0.49\textwidth}
      \begin{block}{Скалярная квантизация}
        \small

        \textbf{Scalar PTQ} --- равномерная сетка
        и независимое округление весов.
        \par\medskip

        \textbf{Q-Diffusion-like} --- калибровка сетки
        по ошибке восстановления выхода блока.
        \par\medskip

        \textbf{PTQ4DM-like} --- timestep-aware калибровка
        и reconstruction objective.
        \par\medskip

        \textbf{BitsFusion-lite} --- разная битность слоёв
        при заданном среднем бюджете.
      \end{block}
    \end{column}

    \begin{column}{0.47\textwidth}
      \begin{block}{Векторная квантизация}
        \small

        \textbf{Vanilla VQ} --- один кодбук;
        быстрый базовый вариант.
        \par\medskip

        \textbf{AQUATIC-Diff-like} --- аддитивное
        представление несколькими кодбуками
        и калибровка по активациям.
        \par\medskip

        \textbf{VQDM} --- несколько кодбуков,
        activation-aware функционал
        и совместный beam search.
      \end{block}
    \end{column}

  \end{columns}

  \vfill

  \keypoint{
    Like/lite --- адаптации к единому no-fine-tuning протоколу,
    а не полные репродукции исходных методов.
  }
\end{frame}

% ===================================================
% 4. ~60 секунд.
% Главное различие — поиск индексов,
% а не форма представления.
% ===================================================

\begin{frame}[t]{От совместного поиска к остаточной квантизации}
  Общее представление вектора
  $w\in\mathbb{R}^{d}$ --- сумма кодовых слов:

  \[
    \widehat w = \sum_{m=1}^{M} C_m[z_m].
  \]

  \small
  $C_m$ --- кодбук стадии $m$;
  $z_m$ --- индекс выбранного вектора.
  \normalsize

  \begin{columns}[T,onlytextwidth]

    \begin{column}{0.47\textwidth}
      \begin{block}{VQDM: совместный выбор}
        Комбинация $(z_1,\ldots,z_M)$ выбирается
        через \textbf{beam search}.
        \par\smallskip

        Точный, но дорогой поиск при росте
        числа кодбуков и слоёв.
      \end{block}
    \end{column}

    \begin{column}{0.49\textwidth}
      \begin{block}{RVQ-AW: последовательный выбор}
        Каждая стадия квантует \textbf{остаток}:

        \[
          r_0=w,
          \qquad
          r_m=r_{m-1}-C_m[z_m].
        \]
      \end{block}
    \end{column}

  \end{columns}

  \vfill

  \keypoint{
    Компромисс: потенциальная потеря точности
    из-за жадного выбора в обмен на более дешёвое построение.
  }
\end{frame}

% ===================================================
% 5. ~70 секунд.
% Объяснить, почему одинаковые ошибки весов не равноценны.
% ===================================================

\begin{frame}[t]{Как RVQ учитывает активации}
  $X$ --- матрица калибровочных входов слоя;
  $\Delta W=W-\widehat W$;
  $G=X^{\top}X$.

  \[
    \mathcal{L}(\widehat W)
    =
    \left\|X\Delta W^{\top}\right\|_F^2
    =
    \operatorname{tr}\!
    \left(
      \Delta W\,G\,\Delta W^{\top}
    \right).
  \]

  \keypoint{
    Минимизируется ошибка выхода слоя,
    а не только расстояние между исходными
    и квантизованными весами.
  }

  \vspace{1mm}

  \begin{columns}[T,onlytextwidth]

    \begin{column}{0.29\textwidth}
      \begin{block}{Разные размеры $K_m$}
        \small
        Ёмкость кодбуков
        перераспределяется между стадиями
        при близком общем битовом бюджете.
        \[
          b_{\mathrm{index}}
          =\frac{1}{d}\sum_m\log_2 K_m.
        \]
      \end{block}
    \end{column}

    \begin{column}{0.37\textwidth}
      \begin{block}{Full-$G$ recoding}
        \small
        После построения кодбуки фиксируются.
        Coordinate descent уточняет \textbf{индексы}
        по полной $G$, учитывая взаимодействия
        между короткими блоками.
      \end{block}
    \end{column}

    \begin{column}{0.29\textwidth}
      \begin{block}{Hadamard}
        \small
        Ортогональное перемешивание
        координат перед RVQ.
        Применяется не безусловно:
        эффект зависит от конфигурации.
      \end{block}
    \end{column}

  \end{columns}

  \vfill

  {\small\color{Muted}
    Recoding не является fine-tuning:
    на этом шаге меняются только дискретные индексы.
  }
\end{frame}

% ===================================================
% 6. ~50 секунд.
% Сходство с исходной моделью и соответствие тексту
% — разные критерии.
% ===================================================

\begin{frame}[t]{Экспериментальный протокол и метрики}
  \vspace{-2mm}

  \begin{columns}[T,onlytextwidth]

    \begin{column}{0.48\textwidth}
      \begin{block}{Одинаковые условия}
        \small

        \textbf{96 изображений на вариант};
        одинаковые prompts, seed и сэмплирование.
        \par\smallskip

        Калибровка: фиксированные запросы
        и шаги диффузии.
        \par\smallskip

        Скалярные методы:
        \textbf{2, 3, 4 бита}.
        Векторные целевые профили:
        \textbf{2.15, 3.15, 4.15 bpp}.
        \par\smallskip

        \textbf{bpp} учитывает индексы и кодбуки
        в расчёте на квантуемый вес.
        \par\smallskip

        $T_q$ --- время квантизации с калибровкой,
        \textbf{не время инференса}.
      \end{block}
    \end{column}

    \begin{column}{0.48\textwidth}
      \begin{block}{Как оценивается результат}
        \small

        \textbf{CLIP I-I $\uparrow$} ---
        семантическое сходство
        с результатом исходной модели.
        \par\smallskip

        \textbf{LPIPS $\downarrow$} ---
        перцептивное расстояние
        до того же результата.
        \par\smallskip

        \textbf{CLIP T-I $\uparrow$} ---
        соответствие изображения тексту.
        \par\smallskip

        \textbf{PSNR, SSIM $\uparrow$} ---
        пиксельное и структурное сходство.
      \end{block}

      \[
        \score
        =
        \text{CLIP I-I}
        -
        \mathrm{LPIPS}.
      \]
    \end{column}

  \end{columns}

  \vfill

  \keypoint{
    Score используется для компактного ранжирования;
    выводы проверяются по совокупности метрик.
  }
\end{frame}

% ===================================================
% 7. ~75 секунд.
% Сначала качество, затем время.
% Быстрые VQ-бейзлайны не скрывать.
%
% Score скопирован из отчёта;
% не пересчитывается по округлённым CLIP/LPIPS.
%
% Отношения Tq(VQDM)/Tq(RVQ-AW): 3.07, 2.87, 2.37.
% Это wall-clock разных запусков, а не контролируемый hardware benchmark.
% ===================================================

\begin{frame}[t]{Результат: RVQ-AW ближе всех к VQDM}
  {\small
    18 слоёв \midblock;
    три лучших метода по Score для каждого профиля.
  }

  \vspace{1mm}

  \begin{center}
    \small
    \renewcommand{\arraystretch}{1.03}
    \setlength{\tabcolsep}{5mm}

    \begin{tabular}{clrrr}
      \toprule

      \textbf{Профиль}
        & \textbf{Метод}
        & \textbf{bpp}
        & \textbf{Score $\uparrow$}
        & $\boldsymbol{T_q}$\textbf{, с $\downarrow$} \\

      \midrule

      $\approx 2$
        & VQDM
        & 2.02
        & \textbf{0.586}
        & 1857.1 \\

      \rowcolor{Tint}
        & \rvq
        & 2.03
        & 0.564
        & 604.2 \\

        & BitsFusion-lite
        & 2.05
        & 0.468
        & 88.4 \\

      \midrule

      $\approx 3$
        & VQDM
        & 3.02
        & \textbf{0.721}
        & 2846.8 \\

      \rowcolor{Tint}
        & \rvq
        & 3.02
        & 0.704
        & 992.8 \\

        & Vanilla VQ
        & 3.01
        & 0.617
        & 1.91 \\

      \midrule

      $\approx 4$
        & VQDM
        & 4.03
        & \textbf{0.825}
        & 3314.0 \\

      \rowcolor{Tint}
        & \rvq
        & 4.03
        & 0.786
        & 1400.0 \\

        & Vanilla VQ
        & 4.00
        & 0.737
        & 1.57 \\

      \bottomrule
    \end{tabular}
  \end{center}

  \vfill

  \keypoint{
    RVQ-AW --- второе место по Score во всех трёх профилях.
    \newline
    Разрыв с VQDM --- $0.017$--$0.039$ Score.
    Зафиксированный wall-clock ниже в $2.4$--$3.1$ раза,
    но аппаратные конфигурации запусков различались.
  }
\end{frame}

% ===================================================
% 8. ~60 секунд.
% Показать, что проверялись и отрицательные гипотезы.
% ===================================================

\begin{frame}[t]{Что показали промежуточные проверки}
  \footnotesize
  \renewcommand{\arraystretch}{1.12}

  \begin{tabularx}{\textwidth}{@{}p{2.45cm}p{2.75cm}Y@{}}
    \toprule

    \textbf{Проверка}
      & \textbf{Постановка}
      & \textbf{Наблюдение} \\

    \midrule

    \textbf{Разные $K_m$}
      & 10 Linear, noH, $d=8$
      & При почти равном bpp Score вырос
        $0.5682\!\to\!0.6100$ для профиля 2.15
        и $0.7076\!\to\!0.7236$ для 3.15. \\

    \addlinespace[1.1mm]

    \textbf{Full-$G$ recoding}
      & Быстрый совмещённый sweep
      & Внутренний функционал уменьшился в среднем
        на \textbf{47.53\%}.
        Внешние метрики менялись неоднозначно;
        нужен изолированный повтор. \\

    \addlinespace[1.1mm]

    \textbf{Hadamard}
      & 10 Linear, $K=256$, без recoding
      & H16 дал $+0.014$ Score при 3.15 bpp;
        при 2.15 и 4.15 лучшим был noH. \\

    \addlinespace[1.1mm]

    \textbf{Размерность $d$}
      & Старый совмещённый pipeline
      & При $K=256$ вариант $d=16$ проиграл $d=8$
        по качеству и потребовал почти вдвое
        больше времени. \\

    \addlinespace[1.1mm]

    \textbf{Local recode}
      & 10 Linear, совмещённый перебор
      & Уточнение между стадиями выиграло
        лишь \textbf{2 из 20} парных сравнений;
        среднее $\Delta\score=-0.0408$. \\

    \bottomrule
  \end{tabularx}

  \vfill

  \keypoint{
    Главный положительный сигнал --- разные $K_m$ по стадиям.
    Текущий кандидат: $d=8$, stage-wise $K_m$, local recode отключён;
    Hadamard и full-$G$ recoding требуют факторной перепроверки.
  }
\end{frame}

% ===================================================
% 9. ~50 секунд.
% Закончить результатом и повесткой дальнейшей работы.
% ===================================================

\begin{frame}[t]{Текущий вклад и следующий этап}
  \keypoint{
    Последовательная RVQ сохраняет большую часть
    преимущества VQDM по Score
    без совместного beam-search поиска.
  }

  \vspace{2mm}

  \begin{columns}[T,onlytextwidth]

    \begin{column}{0.47\textwidth}
      \begin{block}{Что сделано}
        \small

        Единый \textbf{no-fine-tuning протокол}
        для скалярных и векторных методов.
        \par\medskip

        Кандидат \textbf{RVQ-AW}:
        activation-aware остаточная квантизация,
        распределение ёмкости стадий и recoding.
        \par\medskip

        Сравнение при трёх битовых бюджетах
        и проверки отдельных гипотез.
      \end{block}
    \end{column}

    \begin{column}{0.49\textwidth}
      \begin{block}{Что проверить дальше}
        \small

        Контролируемый совместный перебор
        \textbf{$K$, $d$, Hadamard и recoding}.
        \par\medskip

        Развитие timestep-aware калибровки
        и переход от \midblock{}
        к \textbf{полному U-Net}.
        \par\medskip

        FID/clean-FID на большем числе изображений,
        preference-метрика, память и latency.
        Fine-tuning и distillation ---
        отдельная линия экспериментов.
      \end{block}
    \end{column}

  \end{columns}

  \vfill

  {\small\color{Muted}
    Область текущих выводов:
    \midblock{} и 96 изображений на вариант.
    $T_q$ не нормирован по GPU; ускорение инференса из него не следует.
  }
\end{frame}

% ===================================================
% Подробные метрики по каждому профилю.
% Числа перенесены из сводки.
% ===================================================

\begin{frame}[t]{Метрики по каждому профилю}
  \begin{center}
    \small
    \renewcommand{\arraystretch}{1.15}
    \setlength{\tabcolsep}{2.5mm}

    \begin{tabular}{clrrrrr}
      \toprule

      \textbf{Профиль}
        & \textbf{Метод}
        & \textbf{bpp}
        & \textbf{CLIP I-I}
        & \textbf{CLIP T-I}
        & \textbf{LPIPS}
        & \textbf{Score} \\

      \midrule

      $\approx 2$
        & VQDM
        & 2.02 & 0.831 & 0.305 & 0.244 & 0.586 \\

      \rowcolor{Tint}
        & \rvq
        & 2.03 & 0.819 & 0.305 & 0.254 & 0.564 \\

        & BitsFusion-lite
        & 2.05 & 0.768 & 0.312 & 0.300 & 0.468 \\

      \midrule

      $\approx 3$
        & VQDM
        & 3.02 & 0.885 & 0.308 & 0.164 & 0.721 \\

      \rowcolor{Tint}
        & \rvq
        & 3.02 & 0.879 & 0.308 & 0.175 & 0.704 \\

        & Vanilla VQ
        & 3.01 & 0.828 & 0.319 & 0.212 & 0.617 \\

      \midrule

      $\approx 4$
        & VQDM
        & 4.03 & 0.926 & 0.310 & 0.101 & 0.825 \\

      \rowcolor{Tint}
        & \rvq
        & 4.03 & 0.911 & 0.311 & 0.125 & 0.786 \\

        & Vanilla VQ
        & 4.00 & 0.881 & 0.317 & 0.144 & 0.737 \\

      \bottomrule
    \end{tabular}
  \end{center}

  \vfill

  \keypoint{
    Лидерство по Score не означает лидерства по CLIP T-I:
    сходство с исходной моделью и соответствие тексту ---
    разные критерии.
  }
\end{frame}

% ===================================================
% Битовый бюджет и итоговые профили.
%
% Не смешивать отдельную Hadamard-абляцию
% с итоговыми профилями:
% в источнике они дают разные варианты.
% Причина расхождения не уточнена.
% ===================================================

\begin{frame}[t]{Битовый бюджет и итоговые профили RVQ-AW}
  Стоимость индексов для вектора из $d$ весов:

  \[
    b_{\mathrm{index}}
    =
    \frac{1}{d}
    \sum_{m=1}^{M}\log_2 K_m.
  \]

  {\small
    Фактический bpp дополнительно
    учитывает хранение кодбуков.
  }

  \vspace{2mm}

  \begin{center}
    \small
    \renewcommand{\arraystretch}{1.3}
    \setlength{\tabcolsep}{3mm}

    \begin{tabular}{cclcr}
      \toprule

      \textbf{Целевой bpp}
        & \textbf{Hadamard}
        & \boldmath$(K_1,\ldots,K_M)$
        & \boldmath$d$
        & \textbf{Факт. bpp} \\

      \midrule

      2.15
        & noH
        & $[64,1024]$
        & 8
        & 2.03 \\

      3.15
        & H32
        & $[128,256,512]$
        & 8
        & 3.02 \\

      4.15
        & H16
        & $[512,512,128,128]$
        & 8
        & 4.03 \\

      \bottomrule
    \end{tabular}
  \end{center}

  \vspace{2mm}

  {\small\color{Muted}
    Чистый тест Hadamard и итоговый выбор проводились
    в разных постановках: fixed $K$ без recoding на 10 Linear
    против stage-wise $K$ с recoding на 18 слоях.
    Поэтому различие оптимальных блоков не является противоречием,
    но взаимодействия нужно проверить в едином факторном эксперименте.
  }
\end{frame}

\end{document}
